\chapter{Aperture}

``Mummy, my phone has got really old. I can't even use the internet on it for studies. Can I get a new phone?'' Shiva asked his mother over the phone, sitting on the roof.

``What is the issue? Isn't it already working fine?'' she replied.

Picking the last grain of wheat left on the floor, Shiva said, ``Yes, it is fine for searching and all, but I can't watch YouTube videos. There I can find teachers who teach various subjects.''

``But how much will it cost?''

``Not much, Mummy. There are some really affordable Android phones.'' He put it in his mouth.

``Okay, I will ask your Mama to look for one. Don't even dare to go to the market yourself. Highways are risky,'' she said.

``Fine, Mummy. But I am grown up now. In a few years I'll have to go out to another city. Why do you all keep treating me as if I am a child?''

``You don't know anything, Shiva. All you have seen are books, but the world is bigger and crueller, and you are not smart enough for that.''

``Leave it. I don't even want to discuss it. If you keep behaving like this, how am I ever going to learn anything?'' he said, looking for another stray grain in the dim light of the distant bulb.

``Everything has a time,'' she replied, a faint smile audible over the call.

``You know I have learnt how to ride the bike. I asked Nana to let me take him to the market on the motorcycle, but he refused.''

She laughed. ``He has ridden his bicycle all his life. Best of luck convincing him onto a motorcycle now.''

``One day...'' he sighed. His gaze locked on the bulb, which now had a small canopy made from a plastic bottle, leaking red where the light passed through it, and yellow everywhere else.

``Okay, Mummy, I have to go for dinner now. Mami is calling me,'' he said quickly.

``Okay, sure. But first tell me, when can you come here for the winter vacations?'' she asked.

``No, Mummy. It's not possible. Three of the ten days are already over. Going and coming back takes a couple of days on its own—there will be nothing left to stay with you.''

``It's okay. Have your dinner now,'' she said.

``Okay, Mummy. Good night,'' he said as he walked toward the stairs.

``Good night, son.''

After a few days, Shiva was fixing the brakes of his bicycle while Nani and Nana were sitting in the courtyard in the fading sunlight of the afternoon.

Nana, sitting on a chair, was reading his newspaper. Turning a page, he asked, ``Why don't you give it for repair? The bicycle shop near your school—its owner was my student. He is really good, and I have always got mine repaired by him whenever I needed.''

``But, Nana, this is a geared bicycle. They don't get it. They don't have parts for it,'' he said, removing the rear axle from the dropout.

``Yes, and you have done a diploma in cycle repairing. Who can talk sense to this generation?'' He turned one more page.

``Why do you underestimate my grandson like that? Why can't he fix his cycle?'' Nani intervened, pausing the massage of her knees on the cot.

He pulled the newspaper down and replied, ``Oho... I am underestimating your grandson. I just meant to reduce the work he has to do every few days. Why doesn't he just get it done by some professional?''

``It's his cycle, let him do as he likes,'' Nani said, looking towards Shiva.

``Fine. You can also join him,'' Nana laughed, raising the newspaper again.

The sound of a motorcycle came from the gate before Mama appeared, riding it inside. He parked it beside the veranda, near Shiva, then came back and sat down beside Nani.

``What broke now?'' Mama asked.

``My brake was loose, so I tightened it by moving the wheel slightly back, and now my gears are jammed. This brake design is so bad it can't even be fixed properly,'' Shiva said, housing the wheel back in the dropout.

``I told you to take a simpler one, but you chose this one anyway,'' Mama said, putting a bag down beside Nani.

``It's alright, I will get it done. I have to go to school tomorrow,'' Shiva said, now tightening the nuts.

``Sure you will. But won't you see what I have brought for you?''

Shiva looked back to see a black polythene on the cot. ``What's that?'' he asked.

``See for yourself. But first wash your dirty hands,'' Mama said.

``It's almost done.'' Shiva got up, pulled the rear wheel up by the carrier, spun the pedal, and applied the brake hard. The wheel stopped, but not immediately. Shiva wheeled the cycle from the courtyard over beside the motorcycle, as if that were the expected outcome of his hour-long fix. He returned to them after washing his hands.

He opened the bag and found a phone box. ``Thanks, Mama...'' he said, his voice filled with joy.

``It has more internet. One GB of it for a day,'' Mama said with a smile.

``What does it mean to be more? Can he see different answers in there now?'' Nani asked while pouring more oil into her hand.

``No. The answers will be the same, but he can search more questions now,'' Mama said.

``In our times, books were the only thing. I had to save three months of payment I got from working at a shop as a Munshi to get my MA books. In those times even paisas had value. Now you can get all books on one small screen,'' Nana added, turning one more page, now nearing the end of the whole newspaper.

In the evening, he went up to the roof as usual with a cup of tea in his hand, Rahul joining him this time. He brought a carpet from the bulkhead, which Rahul laid out flat.

``I will bring the pakora. At least we will have something warm here,'' Rahul said, putting down his tea.

Shiva dialed his mother and started his conversation.

``Thanks, Mummy,'' he said with a smile on his face.

``It's okay, son,'' she replied, and continued, ``But how big is it? Is it larger than my phone?''

``Yes, Mummy, it is larger, and it doesn't have any keypad like the one you have. We have to use the screen to write anything.''

``Oh sure. But is it suitable for the studies that you got it for?'' she asked.

``Yes, it is.''

``But don't indulge in playing games on it. Focus on your studies.'' Her warning was stern and clear.

Rahul arrived with the pakora and placed the plate beside Shiva before sitting down.

``No. I don't play games,'' Shiva said.

``Don't tell me. You deducted money from my phone so many times when you were young, just looking for games.''

Shiva took a bite of a pakora. ``Mummy, I was younger at that time.''

``Okay, sure, I know you. Take care. I have to start cooking dinner now—I have already prepared the vegetables,'' she said.

Shiva put the phone down to take a sip of tea while Rahul picked it up and asked, ``Bhaiya, how to download games in it, the shooting one?''
